\chapter{Estructura atómica}

En la naturaleza, la materia se encuentra en una gran variedad de formas. Desde el punto de vista de la química, las sustancias puras se distinguen según sean compuestos o elementos. Los elementos son sustancias que no pueden descomponerse más mediante métodos químicos.

La forma mínima en que se puede encontrar un compuesto se llama molécula y la forma mínima en que se puede encontrar un elemento se llama átomo.

Aunque etimológicamente la palabra átomo significa «sin corte», en el presente siglo se ha establecido que los átomos están constituidos por electrones (partículas con carga eléctrica negativa que giran en órbitas específicas alrededor del núcleo), y un núcleo el cual está a su vez formado por protones (partículas con carga eléctrica positiva) y neutrones (partículas con carga eléctrica neutra).

Los átomos son generalmente neutros porque el número de electrones que giran alrededor del núcleo con sus cargas negativas es igual al número de protones en el núcleo con sus cargas positivas.

El número de protones determina la identidad del elemento químico. Cuando un mismo elemento tiene diferentes números de neutrones, en estos casos se habla de isótopos del elemento.

A cada elemento se le asocia con el número de orden que le corresponde dentro de la tabla periódica. A este número se le conoce como número atómico, que corresponde al número de protones que posee el núcleo, y se representa con la letra $Z$.

A la suma de protones y neutrones se le llama número de masa y se representa con la letra $A$.

Para diferenciar un isótopo de otro se utiliza la notación siguiente:

\begin{center}
\fbox{\begin{minipage}{0.82\textwidth}
  \centering
  \vspace{4pt}
  {\Huge $\displaystyle{}^{A}_{Z}X_{N}$}\par
  \vspace{8pt}
  \begin{tabular}{@{}>{$}c<{$}p{0.66\textwidth}@{}}
    X & símbolo químico;\\[2pt]
    A & número másico: suma de protones y neutrones;\\[2pt]
    Z & número atómico;\\[2pt]
    N & número de neutrones.
  \end{tabular}
  \vspace{5pt}
\end{minipage}}
\end{center}

Por ejemplo, para el hidrógeno-2:
\begin{center}
  {\huge $\displaystyle{}^{2}_{1}\mathrm{H}_{1}$}
  \qquad $Z=1,\quad N=1,\quad A=2,\quad X=\mathrm{H}$.
\end{center}
